\dish{姜醋番茄}
\hiragana{jiāng cù fān qié}
\altdish{Tomato with black vinegar and ginger}
\serves{2}
\prep{10 minutes}
\source{openai-chatgpt}
\label{jiangcu-fanqie}

\begin{ingredients}
  \ingr{2--3}{}{ripe tomatoes, cut into wedges}
  \ingr{1}{\tbsp}{Chinkiang black vinegar}
  \ingr{1}{tsp}{light soy sauce}
  \ingr{\fracH}{tsp}{sugar}
  \ingr{1}{tsp}{ginger, finely grated or julienned}
  \ingr{1}{}{spring onion, finely sliced}
  \ingr{\fracQ}{tsp}{green Sichuan pepper oil, or a pinch ground green Sichuan pepper}
  \ingr{}{}{chili oil (optional)}
  \ingr{}{}{coriander (optional)}
\end{ingredients}


\begin{recipe}
  \begin{enumerate}

  \item Mix vinegar, soy sauce, sugar, ginger, spring onion and
    Sichuan pepper oil.

  \item Toss with tomatoes 10--20~minutes before eating (not
    earlier: salt and soy draw water from the tomatoes).

  \end{enumerate}
\end{recipe}

\notes{} Green Sichuan pepper's brighter, citrusy character
contrasts well with the rich sesame noodles. Red Sichuan pepper
works too, but is warmer and woodier.

Serve alongside \hyperref[majiang-liangmian]{麻酱凉面}.

%\accord{}
